\chapter{简介}

竞赛编程结合了两个主题：
（1）算法的设计和（2）算法的实现。

\key{算法的设计}包括问题求解
和数学思维。
需要具备分析问题和创造性地
解决问题的能力。
解决问题的算法
必须既正确又高效，
而问题的核心往往是
要设计一个高效的算法。

算法的理论知识
对竞赛程序员来说很重要。
通常，一个问题的解决方案是
已知技术与新见解的结合。
竞赛编程中出现的技术
也构成了
算法科学研究的基础。

\key{算法的实现}需要良好的
编程技能。
在竞赛编程中，解决方案
是通过对已实现算法
进行一组测试用例的测试来评分的。
因此，仅有正确的算法思想是不够的，
实现也必须正确。

比赛中良好的编码风格是
简洁明了的。
程序应该快速编写，
因为可用时间不多。
与传统软件工程不同，
程序很短（通常最多几百行代码），
而且赛后不需要维护。

\section{编程语言}

\index{编程语言}

目前，比赛中最流行的
编程语言是 C++、Python 和 Java。
例如，在 Google Code Jam 2017 中，
排名前 3000 的参赛者中，
79\% 使用 C++，
16\% 使用 Python，
8\% 使用 Java \cite{goo17}。
部分参赛者还使用了多种语言。

许多人认为 C++ 是竞赛程序员
的最佳选择，
而且 C++ 几乎在所有比赛系统中都可用。
使用 C++ 的好处是
它是一种非常高效的语言，
而且它的标准库包含
大量数据结构和算法。

另一方面，掌握多种语言
并理解它们各自的优势
也是有益的。
例如，如果问题需要大整数运算，
Python 可能是一个好选择，因为它
内置了大整数运算功能。
尽管如此，大多数编程竞赛中的问题
设置使得
使用特定编程语言
不会带来不公平的优势。

本书所有示例程序都用 C++ 编写，
并经常使用标准库中的
数据结构和算法。
程序遵循 C++11 标准，
该标准在目前大多数比赛中都可用。
如果你还不会 C++ 编程，
现在是开始学习的好时机。

\subsubsection{C++ 代码模板}

竞赛编程中典型的 C++ 代码模板
如下所示：

\begin{lstlisting}
#include <bits/stdc++.h>

using namespace std;

int main() {
    // 在此编写代码
}
\end{lstlisting}

代码开头的 \texttt{\#include} 行
是 \texttt{g++} 编译器的一个特性，
它允许我们包含整个标准库。
因此，不需要单独包含
\texttt{iostream}、\texttt{vector} 和 \texttt{algorithm} 等库，
它们会自动可用。

\texttt{using} 行声明
标准库的类和函数
可以在代码中直接使用。
如果没有 \texttt{using} 行，
例如，我们必须写 \texttt{std::cout}，
但现在只需写 \texttt{cout} 就够了。

可以使用以下命令编译代码：

\begin{lstlisting}
g++ -std=c++11 -O2 -Wall test.cpp -o test
\end{lstlisting}

此命令从源代码 \texttt{test.cpp}
生成二进制文件 \texttt{test}。
编译器遵循 C++11 标准
（\texttt{-std=c++11}），
优化代码（\texttt{-O2}）
并显示可能的错误警告（\texttt{-Wall}）。

\section{输入与输出}

\index{输入与输出}

在大多数比赛中，使用标准流
进行输入和输出。
在 C++ 中，标准流是
用于输入的 \texttt{cin} 和用于输出的 \texttt{cout}。
此外，还可以使用 C 函数
\texttt{scanf} 和 \texttt{printf}。

程序的输入通常由
用空格和换行符分隔的
数字和字符串组成。
可以从 \texttt{cin} 流中
如下读取它们：

\begin{lstlisting}
int a, b;
string x;
cin >> a >> b >> x;
\end{lstlisting}

这种代码总是有效的，
前提是输入中每个元素之间
至少有一个空格或换行符。
例如，上面的代码可以读取
以下两种输入：
\begin{lstlisting}
123 456 monkey
\end{lstlisting}
\begin{lstlisting}
123    456
monkey
\end{lstlisting}
\texttt{cout} 流用于输出，
如下所示：
\begin{lstlisting}
int a = 123, b = 456;
string x = "monkey";
cout << a << " " << b << " " << x << "\n";
\end{lstlisting}

输入和输出有时会成为
程序的瓶颈。
在代码开头添加以下几行
可以使输入和输出更高效：

\begin{lstlisting}
ios::sync_with_stdio(0);
cin.tie(0);
\end{lstlisting}

注意，换行符 \texttt{"\textbackslash n"}
比 \texttt{endl} 更快，
因为 \texttt{endl} 总是会
触发刷新操作。

C 函数 \texttt{scanf}
和 \texttt{printf} 是
C++ 标准流的替代方案。
它们通常稍快一些，
但也更难使用。
以下代码从输入中读取两个整数：
\begin{lstlisting}
int a, b;
scanf("%d %d", &a, &b);
\end{lstlisting}
以下代码输出两个整数：
\begin{lstlisting}
int a = 123, b = 456;
printf("%d %d\n", a, b);
\end{lstlisting}

有时程序需要从输入中读取整行，
可能包含空格。
这可以通过使用
\texttt{getline} 函数来实现：

\begin{lstlisting}
string s;
getline(cin, s);
\end{lstlisting}

如果数据量未知，以下
循环很有用：
\begin{lstlisting}
while (cin >> x) {
    // code
}
\end{lstlisting}
此循环从输入中逐个读取元素，
直到输入中没有
更多数据可用。

在某些比赛系统中，使用文件
进行输入和输出。
一个简单的解决方案是照常
使用标准流编写代码，
但在代码开头添加以下几行：
\begin{lstlisting}
freopen("input.txt", "r", stdin);
freopen("output.txt", "w", stdout);
\end{lstlisting}
之后，程序从文件 ''input.txt'' 读取输入，
并将输出写入文件
''output.txt''。

\section{处理数字}

\index{整数}

\subsubsection{整数}

竞赛编程中最常用的整数类型
是 \texttt{int}，它是一个 32 位类型，
取值范围为 $-2^{31} \ldots 2^{31}-1$，
即约 $-2 \cdot 10^9 \ldots 2 \cdot 10^9$。
如果 \texttt{int} 类型不够用，
可以使用 64 位类型 \texttt{long long}。
其取值范围为 $-2^{63} \ldots 2^{63}-1$，
即约 $-9 \cdot 10^{18} \ldots 9 \cdot 10^{18}$。

以下代码定义了一个
\texttt{long long} 变量：
\begin{lstlisting}
long long x = 123456789123456789LL;
\end{lstlisting}
后缀 \texttt{LL} 表示
该数字的类型为 \texttt{long long}。

使用 \texttt{long long} 类型时
一个常见的错误是
代码中某处仍在使用 \texttt{int} 类型。
例如，以下代码包含
一个微妙的错误：

\begin{lstlisting}
int a = 123456789;
long long b = a*a;
cout << b << "\n"; // -1757895751
\end{lstlisting}

虽然变量 \texttt{b} 的类型是 \texttt{long long}，
但表达式 \texttt{a*a} 中的两个数字
都是 \texttt{int} 类型，结果
也是 \texttt{int} 类型。
因此，变量 \texttt{b} 将
包含错误的结果。
可以通过将 \texttt{a} 的类型
改为 \texttt{long long}，或
将表达式改为 \texttt{(long long)a*a} 来解决问题。

通常比赛问题的设置使得
\texttt{long long} 类型就够用了。
不过，了解 \texttt{g++} 编译器
还提供了 128 位类型 \texttt{\_\_int128\_t}
也是有帮助的，
其取值范围为
$-2^{127} \ldots 2^{127}-1$，即约 $-10^{38} \ldots 10^{38}$。
但并非所有比赛系统都支持此类型。

\subsubsection{模运算}

\index{余数}
\index{模运算}

我们用 $x \bmod m$ 表示
$x$ 除以 $m$ 的余数。
例如，$17 \bmod 5 = 2$，
因为 $17 = 3 \cdot 5 + 2$。

有时，问题的答案是一个
非常大的数，但只需要
输出它"模 $m$"的值，即
答案除以 $m$ 的余数
（例如，"模 $10^9+7$"）。
其思想是，即使实际答案
非常大，
使用 \texttt{int} 和 \texttt{long long} 类型
也就足够了。

余数的一个重要性质是
在加法、减法和乘法中，
可以在运算前取余：

\[
\begin{array}{rcr}
(a+b) \bmod m & = & (a \bmod m + b \bmod m) \bmod m \\
(a-b) \bmod m & = & (a \bmod m - b \bmod m) \bmod m \\
(a \cdot b) \bmod m & = & (a \bmod m \cdot b \bmod m) \bmod m
\end{array}
\]

因此，我们可以在每次运算后取余，
数字永远不会变得太大。

例如，以下代码计算 $n!$，
即 $n$ 的阶乘，模 $m$：
\begin{lstlisting}
long long x = 1;
for (int i = 2; i <= n; i++) {
    x = (x*i)%m;
}
cout << x%m << "\n";
\end{lstlisting}

通常我们希望余数
总是在 $0\ldots m-1$ 之间。
然而，在 C++ 和其他语言中，
负数的余数
要么为零，要么为负。
确保没有负余数的一个简单方法是
先照常计算余数，
如果结果为负则加 $m$：
\begin{lstlisting}
x = x%m;
if (x < 0) x += m;
\end{lstlisting}
但是，这只在代码中有减法
且余数可能为负时才需要。

\subsubsection{浮点数}

\index{浮点数}

竞赛编程中常用的浮点类型
是 64 位的 \texttt{double}
以及 \texttt{g++} 编译器的扩展类型
80 位的 \texttt{long double}。
在大多数情况下，\texttt{double} 就够了，
但 \texttt{long double} 更精确。

答案所需的精度
通常在题目描述中给出。
一个简单的输出方法是使用
\texttt{printf} 函数，
并在格式字符串中指定
小数位数。
例如，以下代码
以 9 位小数
打印 $x$ 的值：

\begin{lstlisting}
printf("%.9f\n", x);
\end{lstlisting}

使用浮点数时的一个困难是
有些数字无法精确地
表示为浮点数，
会出现舍入误差。
例如，以下代码
的结果令人惊讶：

\begin{lstlisting}
double x = 0.3*3+0.1;
printf("%.20f\n", x); // 0.99999999999999988898
\end{lstlisting}

由于舍入误差，
\texttt{x} 的值比 1 稍小，
而正确的值应该是 1。

用 \texttt{==} 运算符
比较浮点数是有风险的，
因为可能两个值本应相等
但由于精度误差而不等。
更好的比较浮点数的方法是
假设当两个数之差
小于 $\varepsilon$ 时
它们就相等，
其中 $\varepsilon$ 是一个很小的数。

在实践中，可以
如下比较数字（$\varepsilon=10^{-9}$）：

\begin{lstlisting}
if (abs(a-b) < 1e-9) {
    // a 和 b 相等
}
\end{lstlisting}

注意，虽然浮点数不精确，
但在一定范围内的整数
仍然可以被精确表示。
例如，使用 \texttt{double}
可以精确表示所有
绝对值不超过 $2^{53}$ 的整数。

\section{缩短代码}

在竞赛编程中，短代码是理想的选择，
因为程序应该尽可能
快地编写。
因此，竞赛程序员经常为数据类型
和代码的其他部分定义更短的名称。

\subsubsection{类型名}
\index{typedef@\texttt{typedef}}
使用 \texttt{typedef} 命令
可以为数据类型指定更短的名称。
例如，\texttt{long long} 这个名字很长，
所以我们可以定义一个更短的名称 \texttt{ll}：
\begin{lstlisting}
typedef long long ll;
\end{lstlisting}
之后，代码
\begin{lstlisting}
long long a = 123456789;
long long b = 987654321;
cout << a*b << "\n";
\end{lstlisting}
可以缩写为：
\begin{lstlisting}
ll a = 123456789;
ll b = 987654321;
cout << a*b << "\n";
\end{lstlisting}

\texttt{typedef} 命令
也可以用于更复杂的类型。
例如，以下代码
为整数向量
指定名称 \texttt{vi}，
为包含两个整数的
\texttt{pair} 指定名称 \texttt{pi}。
\begin{lstlisting}
typedef vector<int> vi;
typedef pair<int,int> pi;
\end{lstlisting}

\subsubsection{宏}
\index{宏}
缩短代码的另一种方法是定义
\key{宏}。
宏意味着代码中的某些字符串
在编译前会被替换。
在 C++ 中，宏使用
\texttt{\#define} 关键字定义。

例如，我们可以定义以下宏：
\begin{lstlisting}
#define F first
#define S second
#define PB push_back
#define MP make_pair
\end{lstlisting}
之后，代码
\begin{lstlisting}
v.push_back(make_pair(y1,x1));
v.push_back(make_pair(y2,x2));
int d = v[i].first+v[i].second;
\end{lstlisting}
可以缩写为：
\begin{lstlisting}
v.PB(MP(y1,x1));
v.PB(MP(y2,x2));
int d = v[i].F+v[i].S;
\end{lstlisting}

宏也可以有参数，
这使得缩短循环和其他
结构成为可能。
例如，我们可以定义以下宏：
\begin{lstlisting}
#define REP(i,a,b) for (int i = a; i <= b; i++)
\end{lstlisting}
之后，代码
\begin{lstlisting}
for (int i = 1; i <= n; i++) {
    search(i);
}
\end{lstlisting}
可以缩写为：
\begin{lstlisting}
REP(i,1,n) {
    search(i);
}
\end{lstlisting}

有时宏会引起难以发现的 bug。
例如，考虑以下计算数字平方的宏：
\begin{lstlisting}
#define SQ(a) a*a
\end{lstlisting}
这个宏\emph{并不}总是按预期工作。
例如，代码
\begin{lstlisting}
cout << SQ(3+3) << "\n";
\end{lstlisting}
对应于代码
\begin{lstlisting}
cout << 3+3*3+3 << "\n"; // 15
\end{lstlisting}

更好的宏版本如下：
\begin{lstlisting}
#define SQ(a) (a)*(a)
\end{lstlisting}
现在代码
\begin{lstlisting}
cout << SQ(3+3) << "\n";
\end{lstlisting}
对应于代码
\begin{lstlisting}
cout << (3+3)*(3+3) << "\n"; // 36
\end{lstlisting}


\section{数学}

数学在竞赛编程中起着重要作用，
如果不具备良好的数学能力，
就不可能成为一名成功的竞赛程序员。
本节讨论一些重要的
数学概念和公式，
这些将在本书后面的章节中使用。

\subsubsection{求和公式}

每个形如
\[\sum_{x=1}^n x^k = 1^k+2^k+3^k+\ldots+n^k,\]
的求和（其中 $k$ 是正整数），
都有一个闭合公式，
是 $k+1$ 次多项式。
例如\footnote{\index{法哈伯公式}
这类求和甚至有一个通用公式，称为\key{法哈伯公式}，
但它太复杂了，这里不做介绍。},
\[\sum_{x=1}^n x = 1+2+3+\ldots+n = \frac{n(n+1)}{2}\]
and
\[\sum_{x=1}^n x^2 = 1^2+2^2+3^2+\ldots+n^2 = \frac{n(n+1)(2n+1)}{6}.\]

\key{等差数列}\index{等差数列}
是指数列中任意两个相邻
数字之差为常数的
数列。
例如，
\[3, 7, 11, 15\]
是公差为 4 的等差数列。
等差数列的和可以用公式
\[\underbrace{a + \cdots + b}_{n \,\, \textrm{个数}} = \frac{n(a+b)}{2}\]
计算，
其中 $a$ 是第一个数，
$b$ 是最后一个数，
$n$ 是数字的个数。
例如，
\[3+7+11+15=\frac{4 \cdot (3+15)}{2} = 36.\]
该公式基于这样一个事实：
和由 $n$ 个数组成，
每个数的平均值为 $(a+b)/2$。

\index{等比数列}
\key{等比数列}是指数列中
任意两个相邻数字之比
为常数的数列。
例如，
\[3,6,12,24\]
是公比为 2 的等比数列。
等比数列的和可以用公式
\[a + ak + ak^2 + \cdots + b = \frac{bk-a}{k-1}\]
计算，
其中 $a$ 是第一个数，
$b$ 是最后一个数，
相邻数字之间的比率为 $k$。
例如，
\[3+6+12+24=\frac{24 \cdot 2 - 3}{2-1} = 45.\]

这个公式的推导如下。令
\[ S = a + ak + ak^2 + \cdots + b .\]
两边乘以 $k$，得到
\[ kS = ak + ak^2 + ak^3 + \cdots + bk,\]
解方程
\[ kS-S = bk-a\]
即得公式。

等比数列求和的一个特例是公式
\[1+2+4+8+\ldots+2^{n-1}=2^n-1.\]

\index{调和级数}

\key{调和级数}是形如
\[ \sum_{x=1}^n \frac{1}{x} = 1+\frac{1}{2}+\frac{1}{3}+\ldots+\frac{1}{n}\]
的求和。

调和级数的一个上界是 $\log_2(n)+1$。
即，我们可以
修改每一项 $1/k$，
使 $k$ 变为不超过 $k$ 的
最近的 2 的幂。
例如，当 $n=6$ 时，可以
如下估计
求和：
\[ 1+\frac{1}{2}+\frac{1}{3}+\frac{1}{4}+\frac{1}{5}+\frac{1}{6} \le
1+\frac{1}{2}+\frac{1}{2}+\frac{1}{4}+\frac{1}{4}+\frac{1}{4}.\]
这个上界由 $\log_2(n)+1$ 部分组成
（$1$、$2 \cdot 1/2$、$4 \cdot 1/4$ 等），
每部分的值最多为 1。

\subsubsection{集合论}

\index{集合论}
\index{集合}
\index{交集}
\index{并集}
\index{差集}
\index{子集}
\index{全集}
\index{补集}

\key{集合}是元素的汇集。
例如，集合
\[X=\{2,4,7\}\]
包含元素 2、4 和 7。
符号 $\emptyset$ 表示空集，
$|S|$ 表示集合 $S$ 的大小，
即集合中元素的个数。
例如，在上面这个集合中，$|X|=3$。

如果集合 $S$ 包含元素 $x$，
我们写 $x \in S$，
否则写 $x \notin S$。
例如，在上面的集合中
\[4 \in X \hspace{10pt}\textrm{且}\hspace{10pt} 5 \notin X.\]

\begin{samepage}
可以通过集合运算构造新集合：
\begin{itemize}
\item \key{交集} $A \cap B$ 由
同时在 $A$ 和 $B$ 中的元素组成。
例如，如果 $A=\{1,2,5\}$，$B=\{2,4\}$，
则 $A \cap B = \{2\}$。
\item \key{并集} $A \cup B$ 由
在 $A$ 或 $B$ 或两者中的元素组成。
例如，如果 $A=\{3,7\}$，$B=\{2,3,8\}$，
则 $A \cup B = \{2,3,7,8\}$。
\item \key{补集} $\bar A$ 由
不在 $A$ 中的元素组成。
补集的解释取决于
\key{全集}，它包含所有可能的元素。
例如，如果 $A=\{1,2,5,7\}$，全集为
$\{1,2,\ldots,10\}$，则 $\bar A = \{3,4,6,8,9,10\}$。
\item \key{差集} $A \setminus B = A \cap \bar B$
由在 $A$ 中但不在 $B$ 中的元素组成。
注意，$B$ 可以包含不在 $A$ 中的元素。
例如，如果 $A=\{2,3,7,8\}$，$B=\{3,5,8\}$，
则 $A \setminus B = \{2,7\}$。
\end{itemize}
\end{samepage}

如果 $A$ 的每个元素也属于 $S$，
我们说 $A$ 是 $S$ 的\key{子集}，
记作 $A \subset S$。
集合 $S$ 总是有 $2^{|S|}$ 个子集，
包括空集。
例如，集合 $\{2,4,7\}$ 的子集是
\begin{center}
$\emptyset$,
$\{2\}$, $\{4\}$, $\{7\}$, $\{2,4\}$, $\{2,7\}$, $\{4,7\}$ and $\{2,4,7\}$.
\end{center}

一些常用的集合有
$\mathbb{N}$（自然数）、
$\mathbb{Z}$（整数）、
$\mathbb{Q}$（有理数）和
$\mathbb{R}$（实数）。
集合 $\mathbb{N}$
可以有两种定义方式，
取决于具体情况：
$\mathbb{N}=\{0,1,2,\ldots\}$
或 $\mathbb{N}=\{1,2,3,...\}$。

我们还可以使用如下规则构造集合
\[\{f(n) : n \in S\},\]
其中 $f(n)$ 是某个函数。
该集合包含所有形如 $f(n)$ 的元素，
其中 $n$ 是 $S$ 中的元素。
例如，集合
\[X=\{2n : n \in \mathbb{Z}\}\]
包含所有偶整数。

\subsubsection{逻辑}

\index{逻辑}
\index{否定}
\index{合取}
\index{析取}
\index{蕴含}
\index{等价}

逻辑表达式的值要么是
\key{真}（1），要么是
\key{假}（0）。
最重要的逻辑运算符是
$\lnot$（\key{否定}）、
$\land$（\key{合取}）、
$\lor$（\key{析取}）、
$\Rightarrow$（\key{蕴含}）和
$\Leftrightarrow$（\key{等价}）。
下表展示了这些运算符的含义：

\begin{center}
\begin{tabular}{rr|rrrrrrr}
$A$ & $B$ & $\lnot A$ & $\lnot B$ & $A \land B$ & $A \lor B$ & $A \Rightarrow B$ & $A \Leftrightarrow B$ \\
\hline
0 & 0 & 1 & 1 & 0 & 0 & 1 & 1 \\
0 & 1 & 1 & 0 & 0 & 1 & 1 & 0 \\
1 & 0 & 0 & 1 & 0 & 1 & 0 & 0 \\
1 & 1 & 0 & 0 & 1 & 1 & 1 & 1 \\
\end{tabular}
\end{center}

表达式 $\lnot A$ 的值与 $A$ 相反。
表达式 $A \land B$ 在 $A$ 和 $B$
都为真时为真，
表达式 $A \lor B$ 在 $A$ 或 $B$
或两者都为真时为真。
表达式 $A \Rightarrow B$ 在
$A$ 为真时 $B$ 也为真的情况下为真。
表达式 $A \Leftrightarrow B$ 在
$A$ 和 $B$ 同时为真或同时为假时为真。

\index{谓词}

\key{谓词}是一个取决于其参数
而为真或假的表达式。
谓词通常用大写字母表示。
例如，我们可以定义谓词 $P(x)$，
当且仅当 $x$ 是素数时为真。
使用这个定义，$P(7)$ 为真，但 $P(8)$ 为假。

\index{量词}

\key{量词}将逻辑表达式
与集合中的元素联系起来。
最重要的量词是
$\forall$（\key{对所有}）和
$\exists$（\key{存在}）。
例如，
\[\forall x (\exists y (y < x))\]
表示对于集合中的每个元素 $x$，
都存在集合中的元素 $y$
使得 $y$ 小于 $x$。
这在整数集中成立，
但在自然数集中不成立。

使用上述记号，
我们可以表达各种逻辑命题。
例如，
\[\forall x ((x>1 \land \lnot P(x)) \Rightarrow (\exists a (\exists b (a > 1 \land b > 1 \land x = ab))))\]
表示如果数 $x$ 大于 1
且不是素数，
则存在大于 $1$ 的数 $a$ 和 $b$
使得它们的乘积为 $x$。
这个命题在整数集中成立。

\subsubsection{函数}

函数 $\lfloor x \rfloor$ 将数字 $x$
向下取整为整数，函数
$\lceil x \rceil$ 将数字 $x$
向上取整为整数。例如，
\[ \lfloor 3/2 \rfloor = 1 \hspace{10pt} \textrm{且} \hspace{10pt} \lceil 3/2 \rceil = 2.\]

函数 $\min(x_1,x_2,\ldots,x_n)$
和 $\max(x_1,x_2,\ldots,x_n)$
分别给出 $x_1,x_2,\ldots,x_n$ 中
的最小值和最大值。
例如，
\[ \min(1,2,3)=1 \hspace{10pt} \textrm{且} \hspace{10pt} \max(1,2,3)=3.\]

\index{阶乘}

\key{阶乘} $n!$ 可以定义为
\[\prod_{x=1}^n x = 1 \cdot 2 \cdot 3 \cdot \ldots \cdot n\]
或递归地定义为
\[
\begin{array}{lcl}
0! & = & 1 \\
n! & = & n \cdot (n-1)! \\
\end{array}
\]

\index{斐波那契数}

\key{斐波那契数列}
%\footnote{Fibonacci (约 1175--1250) 是一位意大利数学家。}
在许多场合中出现。
它们可以递归地定义如下：
\[
\begin{array}{lcl}
f(0) & = & 0 \\
f(1) & = & 1 \\
f(n) & = & f(n-1)+f(n-2) \\
\end{array}
\]
前几个斐波那契数是
\[0, 1, 1, 2, 3, 5, 8, 13, 21, 34, 55, \ldots\]
斐波那契数也有一个
闭合公式，有时被称为
\index{比内公式} \key{比内公式}：
\[f(n)=\frac{(1 + \sqrt{5})^n - (1-\sqrt{5})^n}{2^n \sqrt{5}}.\]

\subsubsection{对数}

\index{对数}

数字 $x$ 的\key{对数}
记为 $\log_k(x)$，其中 $k$ 是
对数的底。
根据定义，
$\log_k(x)=a$ 当且仅当 $k^a=x$。

对数的一个有用性质是
$\log_k(x)$ 等于
将 $x$ 除以 $k$ 直到
得到 1 所需的次数。
例如，$\log_2(32)=5$，
因为需要 5 次除以 2：

\[32 \rightarrow 16 \rightarrow 8 \rightarrow 4 \rightarrow 2 \rightarrow 1 \]

对数常用于算法分析，
因为许多高效算法
在每一步都将某物减半。
因此，我们可以使用对数
来估计此类算法的效率。

乘积的对数是
\[\log_k(ab) = \log_k(a)+\log_k(b),\]
因此，
\[\log_k(x^n) = n \cdot \log_k(x).\]
此外，商的对数是
\[\log_k\Big(\frac{a}{b}\Big) = \log_k(a)-\log_k(b).\]
另一个有用的公式是
\[\log_u(x) = \frac{\log_k(x)}{\log_k(u)},\]
利用此公式，如果有某种方法
计算固定底数的对数，
就可以计算任意底数的对数。

\index{自然对数}

数字 $x$ 的\key{自然对数} $\ln(x)$
是底为 $e \approx 2.71828$ 的对数。
对数的另一个性质是
整数 $x$ 在 $b$ 进制下的位数为
$\lfloor \log_b(x)+1 \rfloor$。
例如，$123$ 在 $2$ 进制下
的表示为 1111011，
$\lfloor \log_2(123)+1 \rfloor = 7$。

\section{比赛与资源}

\subsubsection{IOI}

国际信息学奥林匹克竞赛（IOI）
是面向中学生的
年度编程竞赛。
每个国家可以派遣一支由
四名学生组成的代表队。
通常有来自
80 个国家的约 300 名参赛者。

IOI 由两场五小时的竞赛组成。
在两场竞赛中，参赛者需要
解决三个不同难度的算法任务。
任务被分为若干子任务，
每个子任务分配有相应的分数。
虽然参赛者被分成代表队，
但他们以个人身份参赛。

IOI 大纲 \cite{iois} 规定了
IOI 任务中可能出现的主题。
本书涵盖了 IOI 大纲中
几乎所有主题。

IOI 的参赛者通过
国内竞赛选拔。
在 IOI 之前，还会组织许多区域性竞赛，
如波罗的海信息学奥林匹克竞赛（BOI）、
中欧信息学奥林匹克竞赛（CEOI）
和亚太信息学奥林匹克竞赛（APIO）。

一些国家为未来的 IOI 参赛者
组织在线练习赛，
如克罗地亚信息学公开赛 \cite{coci}
和美国计算机奥林匹克竞赛 \cite{usaco}。
此外，还有大量来自波兰竞赛的
题目集可在线获取 \cite{main}。

\subsubsection{ICPC}

国际大学生程序设计竞赛（ICPC）
是面向大学生的
年度编程竞赛。
每支参赛队伍由三名学生组成，
与 IOI 不同，学生们需要合作完成；
每支队伍只有一台计算机。

ICPC 分为若干阶段，最终
最优秀的队伍被邀请参加世界总决赛。
虽然竞赛有数万名参赛者，
但只有少数\footnote{决赛名额的确切数量每年不同；2017 年有 133 个决赛名额。}名额可用，
因此在某些地区，
仅进入决赛就是一项巨大的成就。

在每场 ICPC 竞赛中，队伍有五小时的时间
来解决大约十个算法问题。
只有当解决方案高效地解决了
所有测试用例时，才会被接受。
在竞赛期间，参赛者可以查看其他队伍的结果，
但最后一小时记分牌会被冻结，
无法看到最后一次提交的结果。

ICPC 中可能出现的主题
不像 IOI 那样明确。
无论如何，ICPC 显然
需要更多知识，尤其是更多的数学能力。

\subsubsection{在线竞赛}

还有许多向所有人开放的在线竞赛。
目前最活跃的竞赛网站是 Codeforces，
大约每周组织一次竞赛。
在 Codeforces 中，参赛者被分为两个组别：
初学者参加 Div2，更有经验的程序员参加 Div1。
其他竞赛网站包括 AtCoder、CS Academy、HackerRank 和 Topcoder。

一些公司组织带有线下决赛的在线竞赛。
此类竞赛的例子有 Facebook Hacker Cup、
Google Code Jam 和 Yandex.Algorithm。
当然，公司也利用这些竞赛进行招聘：
在竞赛中表现良好
是证明自己能力的好方式。

\subsubsection{书籍}

除本书外，已经有一些专注于
竞赛编程和算法问题求解的书籍：

\begin{itemize}
\item S. S. Skiena and M. A. Revilla:
\emph{Programming Challenges: The Programming Contest Training Manual} \cite{ski03}
\item S. Halim and F. Halim:
\emph{Competitive Programming 3: The New Lower Bound of Programming Contests} \cite{hal13}
\item K. Diks et al.: \emph{Looking for a Challenge? The Ultimate Problem Set from
the University of Warsaw Programming Competitions} \cite{dik12}
\end{itemize}

前两本书面向初学者，
而最后一本书包含高级内容。

当然，一般算法书籍
也适合竞赛程序员。
一些流行的书籍包括：

\begin{itemize}
\item T. H. Cormen, C. E. Leiserson, R. L. Rivest and C. Stein:
\emph{Introduction to Algorithms} \cite{cor09}
\item J. Kleinberg and É. Tardos:
\emph{Algorithm Design} \cite{kle05}
\item S. S. Skiena:
\emph{The Algorithm Design Manual} \cite{ski08}
\end{itemize}
